\section{Preliminaries}

\begin{itemize}
  \item AI (Artificial Intelligence): The area of computer science that involves the development of machines that have the capability of executing tasks traditionally handled by human intelligence, including reasoning, learning and decision-making.

\vspace{1em}

\item Black Box: A system (usually an AI model) the inner mechanisms of which are found to be opaque or not readily interpretable- one can observe the input and the output, but not the decision-making mechanisms of the system.

\vspace{1em}

\item ML (Machine Learning): A sub-specific section of AI which allows the computer to learn trends of information and perform better in the task without being directly programmed.

\vspace{1em}

\item LLM (Large Language Model): An AI model (such as GPT) uses massive volumes of written material in order to learn its methods of comprehending and producing human-like language to accomplish various similar tasks, such as writing, translating, and answering questions.

\vspace{1em}

\item NLP (Natural Language Processing): It is a software application that can decipher handwriting and translate written text into a voice. One of the fields of AI is giving computers the ability to comprehend, reason, and construct human language.

\vspace{1em}

\item Graph: A representation of relationships formed by defining connection layers between many objects (latents), known as nodes, via a relationship followed between them, called edges.

\vspace{1em}

\item DNN (Deep Neural Network): An artificial intelligence application that naturally learns multifaceted patterns of large datasets in an autonomous manner produces a form of machine intelligence, one that possesses several layers of artificial neurons.

\vspace{1em}

\item Retrieval-Augmented Generation: RAG [3] and dense passage retrieval [8] set the precedent of conditioning generation on
retrieved evidence. Production systems use sparse lexical scoring (BM25 [4]) and dense biencoders [5] indexed for fast search [9] and ranked lists combined with reciprocal rank fusion [6] and reranked by cross-encoders [7] with diversity promoted by maximal marginal relevance [10].
In order to boost recall, hypothetical document embeddings (HyDE) [11] are applied to expand queries. Our
framework uses this well-established retrieval stack but it doesn't end there: it feeds the grounded
Converts each state into a state model and records the origin of each recovered document.

\vspace{1em}

\item World Models and Planning: World models learn how to operate in an environment in a latent space and plan through imagining rollouts [12]. The latent recurrent dynamics support
sample-efficient control. Chain-of-thought prompts [16] and Tree-of-thought reasoning [15]
Introduce intentional multi-step search to LLM. We use a learned clinical latent-dynamics model as a transition function of a tree-of-thought beam search to connect these threads, turning
Converted Reactive Answering to Prospective Reasoning.

\vspace{1em}

\item Uncertainty and Interpretability: Deep ensembles [17] continue to be a powerful, simple method for predictive uncertainty, and contemporary
It is known that networks are miscalibrated if they are not corrected [18]. Shapley values [19,20] provide
an axiomatic credit assignment which we adapt to assign answers to retrieved documents—
The technical underpinning for original source tracking. Self-RAG [21] injects self-reflective critique. Our
Calibrated ensembles and Shapley attribution are integrated as a part of framework, instead of being added as a posthoc add-on.

\vspace{1em}

\item Clinical QA Resources: DDXPlus [22] offers synthetic patients with evidences and differentials on a large scale; and
MedMCQA [23] and MedQA [24] offer exam style questions along with explanations. We use these
as the dynamics training source, the evidence corpus and the held out evaluation
set.


\end{itemize}





\section{Review of Existing Research}


Li et al. examined the ways in which it is possible to establish and sustain data to make AI systems a reliable one \cite{Liang2024Traceability}. They identified that most of the issues in AI stem not from the flaws of the algorithms but rather from biased, partial or poorly self-reported data. The authors devised methods of improving the data lifecycle like data cleaning and adding to, labeling, verification and active learning by DARPA Shapley-based metrics. They argue that AI systems depending on tight validation and good data management are significant in providing safe, reliable and fair AI systems.
 
\vspace{1em} 
George et al. \cite{George2024} investigated the problem of increasing recurrence of the AI accountability, such as the misleading information, transparency of the data, and metadata copyright infringements \cite{George2024}. They argued that the problems negatively criticize justice, trust and moral responsibility in significant ways. These authors suggested clear terms of data management, equitable copyright proposal on excellent works, together with strategies, which should be used to locate counterfeit information. They emphasized that governments, businesses, and universities should collaborate and ensure that AI is developed in a manner that is ethical and fosters responsibility to the world community, equitable compensation to designers, and belief of people in the development of AI technologies.

\vspace{1em} 
Chen et al. \cite{ChenSundar2023} examined the role of quality and performance bias labelling on the extent to which individuals trust AI systems\cite{ChenSundar2023}. In an experiment that was conducted with an outcome of 430 people with a fabricated facial expression classifier, it was found that high-quality labeling was only more convincing to people when the AI was not biased. Performance bias severely reduced this by a big margin with data source cues creating insignificant effects. The authors concluded that it is not only necessary to include clear labeling but the fair performance of building user trust and advancing the trustful AI design.

\vspace{1em} 
Liang et al. examined the inability of AI systems to be tracked and relied on and associated this with non-transparent and non-uniform data pipelines \cite{Liang2022}. They emphasized that although automated model development has been developed and improved, data preparation is performed manually and has concealed errors. Most probably to make everything fair and responsible, authors proposed some systematic and auditable methods to operate the data that involve data gathering, data annotating, and evaluation. Their findings reveal that transparent and controlled data pipelines render AI significantly more trustworthy and predictable, as well as ethical and regulation-wise.

\vspace{1em} 
Suneja et al. examined issues of reliability of AI systems in understanding source code\cite{Suneja2021}. They discussed such issues as noisy datasets, semantic misrepresentations, brittle outputs, as well as bad assessment measures. They proposed a complete roadmap, which would address the issue of data, the models, and the evaluation pipeline by tracing the data sources, model which knows the program structure, the analysis of a program both static and dynamic, and evaluation that specializes in developers. Their structure considers legal and ethical questions and brings it to a firmer, simpler perception and makes it more reliable. The paper provides practical recommendations and a list of various classes of dependability issues to aid the AI-for-code studies and applications.

\vspace{1em} 

Heng et al. examined the challenges associated with the comprehension of the black box models in deep learning and highlighted its insufficiency in such spheres as healthcare and finance \cite{Hassija2024}. They took into account the two major XAI methods and divided them into global and local interpretability. They also made some comparisons between ante-hoc and post-hoc explanations and model-specific and model-agnostic explanations such as LIME, SHAP, and model-distillation. They suggested structures of assisting users in gaining specific knowledge and encouraging justice and trust in smart systems. Ultimately, their main focus was to generate AI, which would be transparent, human-focused, and defensible.

\vspace{1em} 
Kale et al. discussed the problem of the black-box of AI and machine learning when a set of computers undergoes judgment without revealing a clear understanding of how it functions\cite{Kale2023}. They paid attention to the fact that it is difficult to manage prejudice, equity, and responsibility within such systems. They recommended the use of provenance documentation in order to monitor the source of data, process and usage. This would render the data more visible, reproducible and discernible. The literature review and bibliometric analysis conducted by the authors conclude that provenance can render AI more trusted, ethical, and easy to track. They indicated that it must be taken as a best practice to develop smart systems that are simple to grasp and trust.

\vspace{1em} 


Lecue considered how Knowledge Graphs (KGs) may apply to improve AI understanding by using context, causation and reasoning\cite{Lecue2020}. He discussed the incessant challenge of putting AI systems into comprehension by non-experts. He mentioned that the conventional methods of deep learning complicate the issue of trust and accountability among the individual members. The semantic linking of KGs provides a way through which AI systems can provide vocabulary of what is decidedly easier to comprehend. The study by Lecue reveals that critical areas can be more comprehensible, reliable, and trusted by the user when the XAI methods are applied to semantic technology.


\vspace{1em} 
The studied subject included the ease with which it is possible to understand deep neural networks conducted by Rauker et al\cite{Rauker2023}. This is a problem as they are black boxes known to be difficult to rely on, maintain and to blame individuals. They reviewed more than 300 studies and categorized ways of making things more accessible to understand by network parts and training stages. They examined such objects as neurons and subnetworks that are located within the brain. The classification relates interpretability to such areas as neuroscience, robustness, and modularity. Diagnostic and benchmarking tools are the subject of the importance the authors have highlighted in an attempt to make AI systems more transparent, credible and aligned.


\vspace{1em} 

Mohsin et al. researched the security and quality issues in big language model generated code, such as GitHub Copilot and ChatGPT\cite{Mohsin2024}. According to them, testing on open data sets usually creates runaway coding and loopholes. To address this, the authors proposed a safe behavioral learning architecture grounded on ICL patterns that evaluates various LLMs across varying languages. Their study proved that in context learning gives things stability and becomes less vulnerable yet, difficult to use since people lack so much knowledge on security matters.

\vspace{1em} 
By entailing structured knowledge graphs on making huge language models more faithful and easy to understand, Luo et al\cite{Luo2024RoG}. investigated the issue of enhancing the reasoning capabilities of language models. They recommended three-stage pipeline planning, retrieving and reasoning. Pathways of candidate relations are constructed in this pipeline, subgraphs of interest are accessed and only the checked facts are reasoned by the LLMs. Fewer hallucinations are made in this technique, making it possible to follow chains of inferences, and answers are more precise on knowledge-graph benchmarks. The approach enables model distillation, but can support a variety of LLM architectures and commercials enable the deployment of AI reasoning systems in a manner which can be both usable and understandable.

\vspace{1em} 
The authors of Otter investigated the possibility to simplify multimodal large models (LMMs) to be more text, image, and video followers\cite{Li2023Otter}. They developed Otter, which is a Flamingo framework and Perceiver architecture-based one. The MIMIC-IT dataset containing over three million multimodal instruction-response pairs was used to train it. It was shown in experiments that Otter enhances convergence and generalization and allows individuals to memorize various images or videos. The paper demonstrates that multimodal training in context makes reasoning and flexibility significantly more effective, and closer to handling more general purpose visual-language assistants, which would be able to converse in a unified and context-sensitive manner.

\vspace{1em} 
Zhao considered the explainability of large language models (LLMs) and discussed what some hazards they might cause in future assignments and how they are difficult to comprehend\cite{Zhao2024}. These results took the form of the identification of a model of grouping Transformer-based models into various categories in accordance with the method by which they are trained (fine-tuning and prompting), and whether those models provided local or global explanations. An additional point that Zhao discussed was the method of model evaluation, fix bugs, and refining them to improve. The work has scaling, faithfulness and ethics issues. It also provides a means of organized knowledge base to aid explainability research to proceed and bridge amiss that exists between that furnishable by technological LLMs and that trusted by people.


\vspace{1em} 



Atanasova et al. aimed to research post hoc explainability methodologies in text classification to address the existing gap put by the absence of a convenient model of the selection and application of attitudinal and saliency strategies\cite{Atanasova2024}. They developed a framework of a diagnostic evaluation which addressed faithfulness, sensitivity, stability, and consistency with human reasoning. After that, they applied gradient-based, perturbation-based, and occlusion to various models and data sets. The findings indicate that there is no optimal strategy that is superior to all the others in all the parametric ratings, and consistency and human compatibility have trade-offs. The paper provides guidelines and materials used to follow on the selection of explainability tools in NLP applications.

\vspace{1em} 
Sajjad viewed the extent to which NLP models can be conceptualized at the neuron level, which is why an individual neuron stores and organizes linguistic and conceptual knowledge\cite{Sajjad2022}. The paper compared and categorized the methods into five categories, namely visualization, corpus-based, probing, causality, and miscellaneous. This it did by considering their scope, scalability, supervision and causal influence. Evaluations were done on model control, domain adaption and distillation. Findings indicate that linguistic features are both localized and distributed in architectures and indicate numerous rich overlaps. The paper provides some suggestions on how deep NLP models of sensitive applications can be made more readable, comprehensible, and authoritative.

\vspace{1em}
Singhal et al. (2023) examined the clinical knowledge of large language models (LLMs) for their ability to answer medical questions and assist with medical-related tasks\cite{singhal2023llmclinical}. The authors identified that while LLM's are well suited for general language, they are not yet known to be reliable for medical-specific domains. To overcome this, they created a medically fine-tuned language model called Med-PaLM that was trained on medical datasets and aligned using methods to improve medical performance. Medical benchmark assessments and clinical question-answering activities were used to measure the model, and the outputs were evaluated by healthcare professionals based on the following criteria: accuracy, reasoning quality, completeness, and safety. The results indicated that Med-PaLM contains extensive clinical information and that it frequently generates more clinically relevant and accurate answers than previous models. The study also uncovered some instances of false or misleading outputs, however. In conclusion, the study underscores the promising prospects of AI systems in healthcare and the necessity for continued development and security enhancements before their clinical use in real-world settings.

\vspace{1em}
Ziwei Ji et al. studied the issue of hallucination in Natural Language Generation (NLG) systems, which are AI systems that produce false, misleading, or unsupported information\cite{ji2023hallucination}. The research emphasized the impact of these hallucinations on the trustworthiness and reliability of applications like chatbots, machine translation, text summarization, and question-answering systems, especially in sensitive sectors like healthcare, education, and law. The researchers surveyed to gain insights into the causes, types, and effects of hallucination in AI-generated text. They systematically examined and analysed existing work to categorize the various types of hallucination, describing the causes of each, such as training data limitations, model architecture weaknesses and lack of factual knowledge sources. The study also examined current methods for the identification and suppression of hallucinated outputs, including the use of better evaluation metrics, better model designs, and retrieval-based strategies that use external knowledge. The results showed that hallucination is still a great problem in NLG systems, and there is no single solution which is applicable to all systems. The study concluded that enhancing the reliability of evaluation methods and the methods used to fact check AI-generated content are crucial for improving the safety and trustworthiness of such content in real-world applications.

\vspace{1em}
Lewis et al. (2020) have tackled the issue that big neural language models are often not very useful for knowledge-intensive NLP problems, as their knowledge is incorporated in the parameters of the model, which can be difficult to update, and can produce incorrect or outdated answers\cite{lewis2020rag}. The research was carried out to make factual information more accurate and enable models to use external facts without having to undergo costly retraining when information changes. To solve this problem, the authors proposed Retrieval-Augmented Generation (RAG), a framework that combines a neural retriever with a sequence-to-sequence generator. The generator then generates responses based on the documents retrieved from a large external knowledge base (e.g., Wikipedia). Experimental evaluations on the benchmarks of open-domain question answering, fact verification and other related tasks highlighted the superior performance of RAG over the previous state of the art. The model produced more accurate, informative and factually correct responses, while retaining the flexibility of the neural text generation, demonstrating the effective combination of retrieval mechanisms with language generation.  

\vspace{1em}
The problem of efficient reinforcement learning (RL) agent training in complex environments, where learning directly from high-dimensional observations like images, is very demanding in terms of computational resources and data, was tackled by Ha and Schmidhuber (2018)\cite{ha2018worldmodels}. This research was performed to investigate if an agent can learn a compact internal model of its environment and use that model to learn effective decision making policies after fewer interactions. The authors have come up with a World Model architecture: A Variational Autoencoder (VAE) to compress visual observations, a recurrent neural network (RNN) to predict future states, and a controller optimized with evolutionary strategies. This allows the agent to acquire an internal model of the environment and even train in a virtual environment created by the model. Experiments on car-racing and Doom-like tasks showed that world-models-based agents excelled when engaging with fewer real environment interactions, which helped them learn more efficiently and enabled them to evolve policies effectively.




\section{Summary of Key Findings}
\begin{enumerate}
    \renewcommand{\labelenumi}{\roman{enumi})} % sets numbering to i), ii),iii), iv), v), vi)
    \item Explainability and Interpretability: Deep learning models are often black boxes and more explicable and transparent techniques like. They may be based on LIMC, SHAP, model distillation, post-hoc analysis, interpretability of a neuron, and others.

    \item Knowledge Structures and Graphs:  Semantic relations were used and context-based reasoning assists in enhancing clarity, consistency, and human. Individual knowledge of AI reactions.
    
    \item Trade-offs XAI:  Human alignment and model stability and scalability .This have been a dilemma since one of the approaches to explainability would be most appropriate to these three areas.

    \item Data Quality and Provenance:  Quality data is offered by reliable AI, which is not biased or lacks documentation. Provenance tracking, cleaning up, labeling, verification, and Shapley-style measures are methods of enhancing the credibility of models.

    \item Fairness, Ethics, and Accountability:  There are some ethical issues which are based on bias, misinformation, and copyright law infringements: ethical. There should be policies, cross-government, academic and industrial collaboration.

    \item Security and Robustness:  It is advisable that AI code and large language models be so shielded against vulnerabilities with the help of safe learning structures, contextually as well as strenuous assessment Learning.

    \item Multimodal and Reasoning Capabilities: The input can be text, image, or video based more advantageous in the learning to reason based on a context sensitive mode of training generalization and after-training learning. The level of the hallucinations is reduced and the reason is brought out better as well as structured pipelines are added to these that also involve change and planning.

    \item General Seasoning:  A full-fledged artificial intelligence ecosystem must be put in place. Interpretable, ethically, accountable, robust, to manage data and also multimodal in order to create transparency, safety, and trustworthiness in society.

\end{enumerate}
\vspace{8em} 




    

The existing literature demonstrates strong theoretical and empirical progress toward explainable, trustworthy AI. The multi‐pronged approach,incorporating model, data, and process transparency,is promising. However, the field still needs advances in making explanations both faithful and accessible; establishing common evaluation standards; and operationalizing provenance and knowledge grounding in real systems under legal, ethical, and organizational constraints.
For authors seeking to build on this foundation, contributions that integrate human usability, standardized evaluations, and domain‐realistic deployments will be particularly valuable.


